% TeX root=../main.tex
% Tabela 25: Pronome interrogativo-indefinido
\begin{table}[!h]
	\begin{adjustwidth}{0cm}{0cm}
		\caption[Pronome interrogativo-indefinido]{Pronome interrogativo-indefinido}\label{tab:proninterr}
		\begin{center}
			\begin{tabular}[c]{lllll}
				\toprule
				                              & \multicolumn{2}{c}{animado} & \multicolumn{2}{c}{inanimado}                                                                 \\
				\midrule
				\emph{sg.nom.}                & \cirilex{kUto}              & \ocsex{kUto}                  & \cirilex{CIto}                 & \ocsex{CIto}                 \\
				\rowcolor{gray!10}\emph{acu.} & \cirilex{kogo}              & \ocsex{kogo}                  & \cirilex{CIto}                 & \ocsex{CIto}                 \\
				\emph{gen.}                   & \cirilex{kogo}              & \ocsex{kogo}                  & \cirilex{Ceso, Cesogo, CIso}   & \ocsex{Ceso, Cesogo, CIso}   \\
				\rowcolor{gray!10}\emph{dat.} & \cirilex{komu}              & \ocsex{komu}                  & \cirilex{Cemu, Cesomu, CIsomu} & \ocsex{Cemu, Cesomu, CIsomu} \\
				\emph{inst.}                  & \cirilex{cEmI}              & \ocsex{cEmI}                  & \cirilex{CimI}                 & \ocsex{CimI}                 \\
				\rowcolor{gray!10}\emph{loc.} & \cirilex{komI}              & \ocsex{komI}                  & \cirilex{CemI, CesomI}         & \ocsex{CemI, CesomI}         \\
				\bottomrule
			\end{tabular}
		\end{center}
	\end{adjustwidth}
\end{table}
